\input{../tex-inputs/latex-leseansicht-vorspann}

\section[Arthur und Olga Schnitzler an Gustav Schwarzkopf, 6. 10. 1911]{L04517 Arthur und Olga Schnitzler an Gustav Schwarzkopf, 6. 10. 1911}
\nopagebreak\mylabel{L04517v}
\rehead{ }\normalsize\beginnumbering\briefempfaengerindex{Schwarzkopf, Gustav@\textsc{Schwarzkopf, Gustav}!zzzSchnitzler, Olga@\emph{von Olga Schnitzler}!1911-10-061@{6. 10. 1911}|(be}\briefempfaengerindex{Schwarzkopf, Gustav@\textsc{Schwarzkopf, Gustav}!zzzSchnitzler, Arthur@\emph{von Arthur Schnitzler}!1911-10-061@{6. 10. 1911}|(be}
\toendnotes[C]{\smallbreak\pagebreak[2]}
\correspDesc{Versand  durch Arthur Schnitzler, Olga Schnitzler am 6. 10. 1911 in Semmering
\newline{}Übermittlung  am 6. 10. 1911 in Semmering
\newline{}Erhalt  durch Gustav Schwarzkopf im Zeitraum [7. 10. 1911 – 11. 10. 1911?] in Wien}\toendnotes[C]{\smallbreak}
\Standort{DLA, A:Schnitzler, HS.1985.1.1897.}
\physDesc{Bildpostkarte, 250 Zeichen
\newline{}Handschrift Arthur Schnitzler: Bleistift, deutsche Kurrent
\newline{}Handschrift Olga Schnitzler: Bleistift, lateinische Kurrent
\newline{}Versand: Stempel: »\nobreak{}\oindex{Semmering@\textbf{Semmering}|pwk}Semmering 1, 6. X. 11, 9\nobreak{}«.  }\toendnotes[C]{\smallbreak}
\pstart
           \noindent{}{\pb}\textcolor{gray}{\textbf{Villa Meran\oindex{Villa Meran@\textbf{Villa Meran}|pw}.}}\hspace*{2.5em}\textcolor{gray}{\textbf{Schneeberg\oindex{Schneeberg@\textbf{Schneeberg}|pw}, 2075 m.}}\hspace*{2.5em}\textcolor{gray}{\textbf{Kurhaus Semmering, N.-Oe\oindex{Kurhaus Semmering@\textbf{Kurhaus Semmering}|pw}.}}\pend
           \pstart{}{\pb}Herrn \textsc{Gustav Schwarzkopf}\pend{}\pstart{}\textsc{Wien I}\oindex{Wien@\textbf{Wien}|pw}\pend{}\pstart{}\textsc{Tiefer Graben 23}\oindex{Wien@\textbf{Wien}!I., Innere Stadt@\textbf{I., Innere Stadt}!Tiefer Graben 23@\textbf{Tiefer Graben 23}|pw}\pend{}{\bigskip}\vspace{1em}
\pstart
           \raggedleft{}6. X. 911\pend
           \vspace{0.5em}
\pstart
           Herzliche Grüße! Ich bin hier ſeit So{\geminationn}tag
      bei Olga\pwindex{Schnitzler, Olga 17.\,1.\,1882 Wien – 13.\,1.\,1970 Lugano@\textsc{Schnitzler, Olga} (17.\,1.\,1882 Wien – 13.\,1.\,1970 Lugano)|pw}, der es ſchon \label{K_L04517-1v}\edtext{viel viel
               beſſer}{\lemma{\textnormal{\emph{viel viel
               besser}}}\Cendnote{\textnormal{Olga Schnitzler\pwindex{Schnitzler, Olga 17.\,1.\,1882 Wien – 13.\,1.\,1970 Lugano@\textsc{Schnitzler, Olga} (17.\,1.\,1882 Wien – 13.\,1.\,1970 Lugano)|pwk} litt an einem Katarrh, vgl. A. S.: \emph{Tagebuch}, 30. 9. 1911.}}}\label{K_L04517-1} geht, fahre So{\geminationn}tag wieder
               hinein. Olga\pwindex{Schnitzler, Olga 17.\,1.\,1882 Wien – 13.\,1.\,1970 Lugano@\textsc{Schnitzler, Olga} (17.\,1.\,1882 Wien – 13.\,1.\,1970 Lugano)|pw} wird zur \label{K_L04517-2v}\edtext{Premiere\eventindex{Burgtheater@\textbf{Burgtheater}!Premiere von Das weite Land, 14.10.1911 [I.]@Premiere von Das weite Land, 14.10.1911 [I.]|pwv}}{\lemma{\textnormal{\emph{Premiere}}}\Cendnote{\textnormal{Premiere von Das weite Land, 14.10.1911 [I.]. In: A. S.: \emph{Kulturveranstaltungen}, \url{https://schnitzler-kultur.acdh.oeaw.ac.at/pmb88534.html}.}}}\label{K_L04517-2}
      wahrſcheinlich auch ſchon in Wien\oindex{Wien@\textbf{Wien}|pw}{ }ſein. Auf Wiederſehn.{\\}Ihr \spacefill\mbox{A.}\pend
           \selectlanguage{ngerman}\vspace{1em}
\pstart
           \noindent{}\raggedleft{}{[}hs. O. Schnitzler:{]} Ihre \spacefill\mbox{O. S.}\pend
           \selectlanguage{ngerman}\endnumbering\briefempfaengerindex{Schwarzkopf, Gustav@\textsc{Schwarzkopf, Gustav}!zzzSchnitzler, Olga@\emph{von Olga Schnitzler}!1911-10-061@{6. 10. 1911}|)be}\briefempfaengerindex{Schwarzkopf, Gustav@\textsc{Schwarzkopf, Gustav}!zzzSchnitzler, Arthur@\emph{von Arthur Schnitzler}!1911-10-061@{6. 10. 1911}|)be}\mylabel{L04517h}
\begin{anhang}
\end{anhang}\newcommand{\dateiname}{L04517}\newcommand{\titel}{Arthur und Olga Schnitzler an Gustav Schwarzkopf, 6. 10. 1911}\newcommand{\editorInnen}{Herausgegeben von Jahnke, SelmaMüller, Martin Anton}\input{../tex-inputs/latex-leseansicht-abspann}
